% =========================================================================
% English translation (generated by AI using Claude Opus 4.8 (Max effort)).
% This file was translated from Hungarian to English by AI.
% DISCLAIMER: Please review against the original Hungarian .tex files, before relying on it!
% SOURCE: 3. Numerikus módszerek.tex
% NOTE: The final "Numerical integration" section embeds an external Hungarian
%       PDF (Numerikus_Integralas.pdf) via \includepdf; that PDF is NOT
%       translated here (it is not a .tex source).
% =========================================================================
\documentclass[margin=0px]{article}

\usepackage{listings}
\usepackage[utf8]{inputenc}
\usepackage{graphicx}
\usepackage{float}
\usepackage[a4paper, margin=0.7in]{geometry}
\usepackage{amsthm}
\usepackage{amssymb}
\usepackage{amsmath}
\usepackage{fancyhdr}
\usepackage{setspace}
\usepackage{pdfpages}

\onehalfspacing

\newenvironment{tetel}[1]{\paragraph{#1 \\}}{}

\pagestyle{fancy}
\lhead{\it{CS BSc Final Exam Topics}}
\rhead{3. Numerical methods}

\title{\textbf{{\Large ELTE IK - Computer Science BSc} \vspace{0.2cm} \\ {\huge Final Exam Topics}} \vspace{0.3cm} \\ 3. Numerical methods}
\author{}
\date{}

\begin{document}
\maketitle

\begin{center}
\fbox{\parbox{0.92\textwidth}{\small \textit{Note: This document was translated from Hungarian to English by AI. Please verify the information against the original Hungarian source before relying on it.}}}
\end{center}

\begin{tetel}{Numerical methods}
    Iterative methods for nonlinear equations: fixed-point iterations, Newton iteration. Interpolation: Lagrange and Newton forms. Least squares method. Numerical integration: interpolation formulas, Newton--Cotes formulas, simple and composite formulas.
\end{tetel}

\section{Iterative methods for nonlinear equations}

So far we have dealt with systems of equations in which every equation was linear. Now we will look for methods
to solve equations of the type $f(x) = 0$, where $f \in \mathbb{R} \to \mathbb{R}$. The essence of the methods
will be that, according to some criterion, we produce a number sequence which, under certain
conditions, converges to the root of the equation.\\

\noindent \textbf{Bolzano theorem}: Let $f \in C[a,b]$ and $f(a)f(b) <0$, that is, the function $f$ is not $0$
at the points $a$ and $b$, and is of opposite sign there. Then $f$ has a root in the interval $(a,b)$, that is,
$\exists x^{*} \in (a,b): f(x^{*}) = 0$.\\

\noindent \textbf{Consequence of the Bolzano theorem}: If, in addition to the conditions of the Bolzano theorem, $f$ is also strictly monotone, then
$x^{*}$ exists uniquely (since $f$ is invertible).\\

\noindent \textbf{Brouwer fixed-point theorem}: Let $f : [a,b] \to [a,b]$ and $f \in C[a,b]$. Then $\exists
    x^{*} \in [a,b]: x^{*} = f(x^{*})$.\\

\noindent \textbf{Theorem}: Let $f : [a,b] \to [a,b]$, $f \in C^{1}[a,b]$ and let $f'$ have constant sign. Then $\exists!
    x^{*} \in [a,b]: x^{*} = f(x^{*})$.\\

\noindent \textbf{Fixed-point theorem for [a,b]}: Let $f : [a,b] \to [a,b]$ be a contraction with contraction coefficient $q$.
Then:

\begin{enumerate}
    \item	$\exists! x^{*} \in [a,b]: x^{*} = f^{x*}$,

    \item	$\forall x_{0} \in [a,b]: x_{k+1} = f(x_{k})$ is convergent and $x^{*} = \lim\limits_{k \to \infty} x_{k}$,

    \item	$| x_{k} - x^{*} | \leq q^{k} | x_{0} - x^{*}| \leq q^{k}(b-a)$.
\end{enumerate}

\noindent \textbf{Order-$p$ convergence}: The convergent sequence $(x_{k})$ ($\lim\limits_{k \to \infty} x_{k} = x^{*}$)
converges with order $p$ if

\begin{displaymath}
    \lim\limits_{k \to \infty} \frac{|x_{k+1} - x^{*}|}{|x_{k}-x^{*}|^{p}} = c > 0
\end{displaymath}

\noindent A few remarks on the above definition:

\begin{enumerate}
    \item	$p$ is unique and $p \geq 1$

    \item	for $p=1$ linear, for $p=2$ quadratic, for $1 < p < 2$ superlinear

    \item	In practice the form $|x_{k+1} - x^{*}| \leq M|x_{k} - x^{*}|^{p}$ is used, which means that
          it converges with order at least $p$.
\end{enumerate}

\noindent \textbf{Theorem}: Suppose that the sequence $(x_{k})$ is convergent,
$x_{k+1} = f(x_{k})$ and $f'(x^{*}) = f''(x^{*}) = ... = f^{(p-1)}(x^{*}) = 0$, but $f^{(p)}(x^{*}) \not = 0$.
Then $(x_{k})$ converges with order $p$.

\subsubsection{Newton's method}

\begin{figure}[H]
    \centering
    \includegraphics[width=0.6\linewidth]{../../3. Numerikus módszerek/img/newton_pelda}
    \caption{The idea of Newton's method.}
    \label{fig:newton_pelda}
\end{figure}

In the figure we start from the leftmost point, 4.12; we take the tangent of the function at this point,
then the root of this tangent will be the next point, and so on. In general,
consider the equation of the tangent of the function $f$ at the point $x_{k}$:

\begin{displaymath}
    y - f(x_{k}) = f'(x_{k}) (x-x_{k})
\end{displaymath}

\noindent As we said, the element $x_{k+1}$ of the iteration is given by the root of the tangent at $x_{k}$:

\begin{displaymath}
    0 - f(x_{k}) = f'(x_{k}) (x_{k+1}-x_{k}) \Rightarrow
    - \frac{f(x_{k})}{f'(x_{k})} = x_{k+1} - x_{k} \Rightarrow
    x_{k+1} = x_{k} - \frac{f(x_{k})}{f'(x_{k})}
\end{displaymath}

The above formula is the formula of Newton's method. It can be noted that the above method is also of the form $x_{k+1} = g(x_{k})$
(fixed-point iteration).

It is important to mention that for $f'(x_{k}) = 0$ the method is not defined. In practice $f'(x_{k}) \approx 0$ is also a problem.
If $x^{*}$ is a multiple root, then $f'(x^{*}) = 0$, that is, near $x^{*}$ the value $f'(x_{k})$ gets closer and closer to $0$, numerically
unstable.\\

\noindent \textbf{Monotone convergence theorem}: Suppose that $f \in C^{2}[a,b]$ and
\begin{enumerate}
    \item	$\exists x^{*} \in [a,b] : f(x^{*}) = 0$, that is, there is a root

    \item	$f'$ and $f''$ have constant sign

    \item	let $x_{0} \in [a,b]$ be such that $f(x_{0}) f''(x_{0})) >0$, that is, $f(x_{0}) \not = 0$, and $f(x_{0})$ and $f''(x_{0})$
          have the same sign
\end{enumerate}

\noindent Then Newton's method started from $x_{0}$ converges monotonically to $x^{*}$.

\noindent \textbf{Local convergence theorem}:	Suppose that $f \in C^{2}[a,b]$ and
\begin{enumerate}
    \item	$\exists x^{*} \in [a,b] : f(x^{*}) = 0$, that is, there is a root

    \item	$f'$ has constant sign

    \item	$0<m_{1} \leq |f'(x)| (x \in [a,b])$ lower bound

    \item	$f''(x) \leq M_{2} (x \in [a,b])$ upper bound, $M:= \frac{M}{2m_{1}}$

    \item	$x_{0} \in [a,b]: |x_{0} - x^{*}| < r = \min \left\{\frac{1}{M}, |x^{*}-a|, |x^{*}-b|\right\}$
\end{enumerate}

\noindent Then Newton's method started from $x_{0}$ converges with second order, with error estimate:

\begin{displaymath}
    |x_{k+1} - x^{*}| \leq M|x_{0} - x^{*}|^{2}
\end{displaymath}
\subsubsection{Chord method (regula falsi)}

The goal is still to solve $f(x) = 0$ on a given interval $[a, b]$.
The essence of the procedure is the following. Initially $x_{0} := a, x_{1} := b$, then we draw the line
formed by these points. Let $x_{2}$ be the root of the chord. If $f(x_{2}) = 0$, then we have found the root. If $f(x_{2}) \not = 0$,
then we continue the search in the interval $[x_{0},x_{2}]$ or $[x_{2},x_{1}]$. If $f(x_{0})f(x_{2}) < 0$, then
we continue in the interval $[x_{0},x_{2}]$; if $f(x_{2}) f(x_{1}) < 0$, then in the interval $[x_{2}, x_{1}]$. And so on.

In general: Let $x_{0} := a, x_{1} := b$ and	$f(a)f(b) < 0$. We approximate the function with the lines passing through the points
$(x_{k},f(x_{k}))$ and $(x_{s},f(x_{s}))$, where for $x_{s}$ we have $f(x_{s})f(x_{k})<0$ and $s$ is the largest such index.
We can determine $x_{k+1}$ as follows:

\begin{displaymath}
    x_{k+1} := x_{k} -
    \frac
    {f(x_{k})}
    {\frac{f(x_{k}) - f(x_{s})}{x_{k}-x_{s}}} =
    x_{k} -
    \frac
    {f(x_{k}) (x_{k} -x_{s}) }
    {f(x_{k}) - f(x_{s})}
\end{displaymath}

\noindent \textbf{Theorem}: Let $f \in C^{2}[a,b]$ and
\begin{enumerate}
    \item	$f(a)f(b)<0$
    \item	$M = \frac{M_{2}}{2m_{1}}$, where $0<m_{1} \leq |f'(x)|$ and $f''(x) \leq M_{2} (x \in (a,b))$
    \item	$M(b-a) < 1$
\end{enumerate}

\noindent Then the chord method is convergent, and its error estimate is:

\begin{displaymath}
    |x_{k+1} - x^{*}| \leq \frac{1}{M}(M|x_{0}-x^{*}|)^{k+1}
\end{displaymath}

\subsubsection{Secant method}
The essence of the secant method is that we approximate $f$ with the line passing through the points $(x_{k},f(x_{k}))$ and $(x_{k-1},f(x_{k-1}))$;
the intersection point of the resulting line with the $x$ axis $(x_{k+1})$ will be the next point. This is
essentially the chord method for $s := k - 1$.

\begin{displaymath}
    x_{k+1} :=
    x_{k} -
    \frac
    {f(x_{k}) (x_{k} -x_{k-1}) }
    {f(x_{k}) - f(x_{k-1})}
\end{displaymath}

\noindent \textbf{Theorem}: If the conditions of the local convergence theorem of Newton's method are satisfied, then the secant method
is convergent with order $p = \frac{1+\sqrt{5}}{2}$, and its error estimate:

\begin{displaymath}
    |x_{k+1} - x^{*}| \leq M|x_{k}-x^{*}||x_{k-1}-x^{*}|
\end{displaymath}

\subsubsection{Multivariate Newton's method}

Now we look for the solutions of $F(x) = 0$, where $F \in \mathbb{R}^{n} \to \mathbb{R}^{n}$. Consider the multivariate Newton's method:

\begin{displaymath}
    x^{(k+1)} := x^{(k)} - [F'(x^{(k)})]^{-1}F(x^{(k)})
\end{displaymath}

\noindent where

\begin{displaymath}
    F(x) = \begin{bmatrix}
        f_{1}(x) \\[0.3em]
        f_{2}(x) \\[0.3em]
        ...      \\[0.3em]
        f_{n}(x)
    \end{bmatrix},
    F'(x) = \begin{bmatrix}
        \partial_{1}f_{1}(x) & \partial_{2}f_{1}(x) & ... \\[0.3em]
        \partial_{1}f_{2}(x) & \partial_{2}f_{2}(x) & ... \\[0.3em]
        ...                  & ...                  & ...
    \end{bmatrix}
\end{displaymath}

\noindent In fact we solve the system of equations $F'(x^{(k)}) (x^{(k+1)} - x^{(k)}) = -F(x^{(k)})$.

\section{Interpolation}

In practice the following problem often arises: that we have to work, on a given interval, with a function
that is mostly unknown (only its values at a few points are known) or very costly to compute.
Then, for example, we can do the following: at a few points we compute the value of the function,
then look for a function that is simpler to compute and that fits
the given points. After this, at any further point of the interval we use the values of the fitted function as
approximations of the values of the original function. Such simpler-to-compute functions are e.g. polynomials (polynomial interpolation).

Example application: When making an animation, we do not want to make every single frame ourselves, but only certain frames,
the so-called key frames. On the intermediate frames we want the computer to compute the positions of the individual objects (for example we want
an object to perform uniform linear motion between two given positions).

\subsection{Polynomial interpolation}

\subsubsection{The polynomial interpolation problem}
Let $n \in \mathbb{N}$ and the distinct numbers $x_{k} \in \mathbb{R}, k=0,1,..,n$, the so-called interpolation base points, be given, as well as
the numbers $f(x_{k}),k=0,1,..,n$, the known function values. We look for the polynomial $p_{n}$ of degree at most $n$ $(p_{n} \in P_{n})$
for which:

\begin{displaymath}
    p_{n}(x_{k}) = f(x_{k}) \ \ k=0,1,...,n
\end{displaymath}

\noindent That is, the sought polynomial takes the given function values at the interpolation base points.\\

\noindent \textbf{Theorem}: The above interpolation problem has a unique solution.\\

\subsubsection{Lagrange interpolation}
Lagrange interpolation gives an explicit formula for computing the interpolation polynomial.\\

\noindent \textbf{Lagrange base polynomials}: For given $n \in \mathbb{N}$ and distinct base points $x_{k} \in \mathbb{R}, k=0,1,...,n$,
the Lagrange base polynomials are defined as follows:

\begin{displaymath}
    l_{k}(x) =
    \frac
    {\displaystyle\prod_{\substack{j=0\\j \not = k}}^{n}(x-x_{j})}
    {\displaystyle\prod_{\substack{j=0\\j \not = k}}^{n}(x_{k}-x_{j})}
\end{displaymath}

\noindent for $k=0,1, ..., n$.\\

\noindent The base polynomials are all of degree $n$, and have the following property:

\begin{displaymath}
    l_{k}(x_{j})=\left\{\begin{array}{lr}
        1 & \text{if} \ j=k        \\
        0 & \text{if} \ j \not = k
    \end{array}
    \right\}
\end{displaymath}

\noindent \textbf{Theorem}: The solution of the interpolation problem is the following polynomial, which we call
the Lagrange form of the interpolation polynomial:

\begin{displaymath}
    L_{n}(x) = \sum_{k=0}^{n}f(x_{k})l_{k}(x)
\end{displaymath}

\noindent In what follows, let $[a,b]$ denote the interval spanned by the base points $x_{0}, x_{1}, ..., x_{n}$.\\

\noindent \textbf{Error of Lagrange interpolation}: If $f \in C^{n+1}[a,b]$, then for $\forall x \in [a,b]$
\begin{displaymath}
    |f(x) - L_{n}(x)| \leq \frac{M_{n+1}}{(n+1)!}|\omega_{n}(x)|
\end{displaymath}

\noindent where $\omega_{n}(x) = \displaystyle\prod_{i=0}^{n}(x-x_{i})$, and $M_{k} = \displaystyle\max_{[a,b]}|f^{(k)}(x)|$.\\

\noindent \textbf{Uniform convergence}: Let $f \in C^{\infty}[a,b]$ and let a sequence of base point systems in the interval $[a,b]$
be given: $x_{k}^{(n)}$, $k=0,1,...,n$, $n=0,1,2,...$. Let $L_{n}$ be the Lagrange interpolation polynomial fitted to the
$x_{0}^{(n)}, ... ,x_{n}^{(n)}$ base point system ($n=0,1,2,...$). Then if $\exists M>0$ such that
$M_{n} \leq M^{n} \ \forall n \in \mathbb{N}$, then the sequence $L_{n}$ converges uniformly to the function $f$.\\

\noindent \textbf{Marcinkiewicz's theorem}: For every $f \in C[a,b]$ there exists a base point system defined as above
such that $\Vert f - L_{n}\Vert_{\infty} \to 0$.\\

\noindent \textbf{Faber's theorem}: For every base point system defined as above there is a function $f \in C[a,b]$
such that $\Vert f - L_{n}\Vert_{\infty} \not \to 0$.\\

\noindent \textbf{Divided difference}: Let the distinct base points $x_{k} \in [a,b], k=0,1,...,n$ be given.
The first-order divided differences of the function $f:[a,b] \to \mathbb{R}$ with respect to the given base point system are

\begin{displaymath}
    f[x_{i},x_{i+1}] = \frac{f(x_{i+1})-f(x_{i})}{x_{i+1} - x_{i}}
\end{displaymath}

\noindent the higher-order divided differences are defined recursively. Suppose that the divided differences of order $k-1$
have already been defined, then the $k$-th order divided differences are the following:

\begin{displaymath}
    f[x_{i},x_{i+1},...,x_{i+k}] = \frac{f[x_{i+1},x_{i+2}, ..., x_{i+k}]-f[x_{i},x_{i+1}, ..., x_{i+k-1})]}{x_{i+k} - x_{i}}
\end{displaymath}

\noindent It can be seen that the $k$-th order divided difference rests on $k+1$ base points.\\

If an interpolation base point system with function values is given, then it is worth arranging the corresponding divided differences according to the
table below, and this arrangement also aids the computation.

\begin{table}[H]
    \begin{tabular}{llllllll}
        $x_{0}$                         & $f(x_{0})$                    &                      &     &                                       &     &   & \\
        $x_{1}$                         & $f(x_{1})$                    & $f[x_{0},x_{1}]$     &     &                                       &     &   & \\
        $x_{2}$                         & $f(x_{2})$                    & $f[x_{1},x_{2}]$     &     &                                       &     &   & \\
        \rotatebox[origin=c]{90}{...}   & \rotatebox[origin=c]{90}{...} &                      &     &                                       &     &   & \\
        $x_{k}$                         & $f(x_{k})$                    & $f[x_{k-1},x_{k}]$   & ... & $f[x_{0}, x_{1}, ..., x_{k}]$         &     &   & \\
        \rotatebox[origin=c]{90}{...}   & \rotatebox[origin=c]{90}{...} &                      &     &                                       &     &   & \\
        $x_{n-1}$                       & $f(x_{n-1})$                  & $f[x_{n-2},x_{n-1}]$ & ... & $f[x_{n-1-k}, x_{n-k}, ..., x_{n-1}]$ & ... &
        $f[x_{0}, x_{1}, ..., x_{n-1}]$ &                                                                                                                \\
        $x_{n} $                        & $f(x_{n})$                    & $f[x_{n-1},x_{n}]$   & ... & $f[x_{n-k}, x_{n-k+1}, ..., x_{n}]$   & ... &
        $f[x_{1}, x_{2}, ..., x_{n}]$   & $f[x_{0}, x_{1}, ..., x_{n}]$                                                                                  \\
    \end{tabular}
\end{table}


\noindent \textbf{The Newton form of the interpolation polynomial}: The interpolation polynomial can be written in the following form:
\begin{displaymath}
    N_{n}(x) = f(x_{0}) + \displaystyle\sum_{k=1}^{n}f[x_{0},x_{1}, ..., x_{k}] \omega_{k-1}(x)
\end{displaymath}

\noindent where $\omega_{j}(x) =  (x-x_{0})(x-x_{1})...(x-x_{j})$. We call this form the Newton form of the interpolation polynomial.

\subsubsection{Hermite interpolation}
We can generalize the previous interpolation problem as follows. Let, in addition to the function values, the derivative values of the function up to some order
also be given at the individual base points. Then we look for a polynomial which, together with its derivatives,
fits the given values, that is:\\

\noindent Let $n,m_{0},m_{1},..., m_{n} \in \mathbb{N}$ and the interpolation base points $x_{j} \in \mathbb{R}, j=0,1,...,n$ be given, as well as
the function and derivative values $f^{(k)}(x_{j}) \ \ k=0,1,...,m_{j}-1, \ \ \ j=0,1,...,n$. Let
$m = \displaystyle\sum_{j=0}^{n}m_{j}$. We look for the polynomial $p_{m-1}$ of degree at most ($m-1$) for which:

\begin{displaymath}
    p^{(k)}_{m-1}(x_{j}) = f^{(k)}(x_{j}) \ \ k=0,1,...,m_{j}-1, \ \ \ j=0,1,...,n
\end{displaymath}

\noindent \textbf{Remarks}:
\begin{enumerate}
    \item	If $m_{j}=2, j=0,1,...,n$, then we call the problem Hermite--Fejér interpolation. Then at every base point
          the function and the first derivative value are given. The sought polynomial is of degree at most ($2n+1$).

    \item	If $m_{j}=1, j=0,1,...,n$, then we get back Lagrange interpolation.
\end{enumerate}

\noindent \textbf{Divided difference for repeated base points}: If $x_{k}$ appears $j$ times:
\begin{displaymath}
    f[x_{k}, ... x_{k}] = \frac{f^{(j)}(x_{k})}{j!}
\end{displaymath}

\noindent \textbf{Theorem}: The Hermite interpolation polynomial exists uniquely.\\

\noindent \textbf{Constructing the Hermite interpolation polynomial}: It can easily be written in the Newton form. All we
have to do is start from the table of base points and function values and produce the table of divided differences.
The only difference now is that we list the base point $x_{j}$ $m_{j}$ times.\\

\noindent \textbf{Error of Hermite interpolation}: If $f \in C^{m}[a,b]$, then for $\forall x \in [a,b]:$

\begin{displaymath}
    |f(x) - H_{m-1}(x)| \leq \frac{M_{m}}{m!}|\Omega_{m}(x)|
\end{displaymath}

\noindent where $\Omega_{m}(x) = (x-x_{0})^{m_{0}}(x-x_{1})^{m_{1}}...(x-x_{n})^{m_{n}}$

\section{Least squares method}

The following problem arises in practical tasks. We measure a first-degree function at certain points, but
because of measurement errors these will not lie on one line. Then we look for a line that, in the following
sense, best fits the given set of measurement points.

Let the measurement points $(x_{i},y_{i}), i=1,2,..,m$ be given. We look for the polynomial $p_{1}(x) = a + bx$ of degree at most
one for which the expression

\begin{displaymath}
    \displaystyle\sum_{i=1}^{m}(y_{i}-p_{1}(x_{i}))^{2} =
    \displaystyle\sum_{i=1}^{m}(y_{i}- a - bx_{i})^{2}
\end{displaymath}

\noindent is minimal. This means that we look for the line for which the sum of the squares of the errors of the function values
is minimal.

The general problem is the following.

Given $m,n \in \mathbb{N}$, where $m >> n$ and the measurement points $(x_{i},y_{i}), i= 1,2,...,m$, where the
base points $x_{i}$ are distinct. We look for the polynomial $p_{n}(x) = a_{0} + a_{1}x + ... a_{n}x^{n}$ of degree at most $n$
for which the expression

\begin{displaymath}
    \displaystyle\sum_{i=1}^{m}(y_{i}-p_{n}(x_{i}))^{2}
\end{displaymath}

\noindent is minimal.

To solve the problem, consider a reformulation of it.

Take the system of equations $p_{n}(x_{i}) = y_{i}, i=1,2,...,m)$. This system is linear in the unknown coefficients $a_{i}$,
and indeed overdetermined; its matrix $A$ is a rectangular Vandermonde matrix $A \in \mathbb{R}^{m \times (n+1)}$,
and the right-hand side vector $b \in \mathbb{R}^{m}$ comes from the function values:

\begin{displaymath}
    \begin{bmatrix}
        1                             & x_{1}                         & x_{1}^{2}                     & ...                           & x_{1}^{n}                     \\[0.3em]
        1                             & x_{2}                         & x_{2}^{2}                     & ...                           & x_{2}^{n}                     \\[0.3em]
        1                             & x_{3}                         & x_{3}^{2}                     & ...                           & x_{3}^{n}                     \\[0.3em]
        \rotatebox[origin=c]{90}{...} & \rotatebox[origin=c]{90}{...} & \rotatebox[origin=c]{90}{...} & \rotatebox[origin=c]{90}{...} & \rotatebox[origin=c]{90}{...} \\[0.3em]
        1                             & x_{m}                         & x_{m}^{2}                     & ...                           & x_{m}^{n}                     \\[0.3em]
    \end{bmatrix}
    \begin{bmatrix}
        a_{0}                         \\[0.3em]
        a_{1}                         \\[0.3em]
        a_{2}                         \\[0.3em]
        \rotatebox[origin=c]{90}{...} \\[0.3em]
        a_{n}
    \end{bmatrix}
    =
    \begin{bmatrix}
        y_{0}                         \\[0.3em]
        y_{1}                         \\[0.3em]
        y_{2}                         \\[0.3em]
        \rotatebox[origin=c]{90}{...} \\[0.3em]
        y_{m}
    \end{bmatrix}
\end{displaymath}

With these notations the expression to be minimized is $\Vert Az - b \Vert_{2}^{2}$, where $z = [a_{0},a_{1}, ..., a_{n}]^{T}$
is the vector of the sought coefficients. The solution of the problem is given by the Gaussian normal equations:

\begin{displaymath}
    A^{T}Az = A^{T}b
\end{displaymath}

\noindent The above system of linear equations has to be solved for $z$.\\

\noindent \textbf{The case n=1}: Then the problem is often also called linear regression. In this case
\begin{displaymath}
    A^{T}A = \begin{bmatrix}
        m                    & \sum_{i=1}^{m} x_{i}     \\[0.3em]
        \sum_{i=1}^{m} x_{i} & \sum_{i=1}^{m} x_{i}^{2}
    \end{bmatrix} \ ,
    A^{T}b = \begin{bmatrix}
        \sum_{i=1}^{m} y_{i} \\[0.3em]
        \sum_{i=1}^{m} x_{i}y_ {i}
    \end{bmatrix} \ ,
    z= \begin{bmatrix}
        b \\[0.3em]
        a
    \end{bmatrix}
\end{displaymath}

\section{Numerical integration}

% NOTE: The pages below come from the external Hungarian PDF and are NOT translated.
\includepdf[pages={2-8}]{../../3. Numerikus módszerek/Numerikus_Integralas.pdf}

\end{document}
